% ============================
% LaTeX Template - Problem Set
% ============================
\documentclass[11pt]{article}

% ---------- Page + layout ----------
\usepackage[margin=1in]{geometry}
\usepackage{setspace}
\setstretch{1.1}
\usepackage{parskip}

% ---------- Math ----------
\usepackage{amsmath, amssymb, amsfonts, amsthm, mathtools}

% ---------- Tables ----------
\usepackage{booktabs}

% ---------- Links ----------
\usepackage[hidelinks]{hyperref}

% ---------- Header / footer ----------
\usepackage{fancyhdr}
\pagestyle{fancy}
\fancyhf{}
\lhead{\textbf{\course}}
\chead{\textbf{\pset}}
\rhead{\textbf{\name}}
\cfoot{\thepage}

% =========================
% Metadata (edit if needed)
% =========================
\newcommand{\name}{Tate Mason}
\newcommand{\course}{Econometrics II}
\newcommand{\pset}{Problem Set 2}
\newcommand{\duedate}{February 16, 2026}

% =========================
% Question / solution macros
% =========================
\newcounter{q}
\renewcommand{\theq}{\arabic{q}}

\newcommand{\question}[1]{%
  \refstepcounter{q}
  \vspace{0.75em}
  \noindent{\Large\textbf{Question \theq.} \normalsize\textbf{#1}}%
  \vspace{0.25em}
  \par
}

\newcounter{subq}[q]
\renewcommand{\thesubq}{\alph{subq}}

\newcommand{\subquestion}[1]{%
  \refstepcounter{subq}
  \vspace{0.35em}
  \noindent\textbf{(\thesubq)} \textbf{#1}\par
}

\newenvironment{solution}{%
  \vspace{0.15em}
  \noindent\textit{Solution. }\ignorespaces
}{\vspace{0.5em}\par}

% =========================
% Easy math macros
% =========================
\newcommand{\R}{\mathbb{R}}
\newcommand{\N}{\mathbb{N}}
\DeclareMathOperator{\E}{\mathbb{E}}
\DeclareMathOperator{\Var}{Var}
\DeclareMathOperator{\Cov}{Cov}
\DeclareMathOperator{\plim}{plim}
\DeclareMathOperator*{\argmin}{argmin}
\DeclareMathOperator*{\argmax}{argmax}
\newcommand{\P}{\mathbb{P}}

\newcommand{\vct}[1]{\mathbf{#1}}
\newcommand{\mtx}[1]{\mathbf{#1}}
\newcommand{\dd}{\,\mathrm{d}}
\newcommand{\eps}{\varepsilon}
\newcommand{\1}{\mathbf{1}}

% =========================
% Document
% =========================
\begin{document}

% ---------- Title block ----------
{\Large\textbf{\course} \hfill \textbf{\pset}}\\
\textbf{Name:} \name \hfill \textbf{Due:} \duedate\\
\rule{\textwidth}{0.4pt}
\question{This exercise will give you practice coding and manipulating discrete choice models in the simplest
setting - a binary choice problem. Here we will examine the determinants of attending a 4-yr college
relative to a 2-yr college among post-secondary attendees. Use “dataHW2 Problem1.mat”. The data
contains an outcome variable Y which is an indicator for attending a 4-yr college. Xi are student
observable characteristics including a constant (these coefficients must vary across choices). The
second column of Xi is an indicator for whether either parent graduated from a 4-yr college, and
the third column of Xi is the student’s high school GPA. Z contains choice specific attributes (the
coefficients are common across choices). The first column of Z captures the distance to the closest
2-year college in miles. The second column of Z captures the distance to the closest 4-year college.}

\subquestion{Estimate a logit of 4-yr college choice using Xi and Z as regressors. Report the coefficient estimates and standard errors.}
\begin{solution}
  \begin{tabular}{lcc}
    \toprule
    Variable & Coefficient & Std. Error \\
    \midrule
    Constant            & -1.1903 & 0.2545 \\
    Parent College Grad & 1.2518  & 0.0799 \\
    High School GPA     & 0.6638  & 0.0903 \\
    Distance            & -0.0226 & 0.040  \\
    \bottomrule
  \end{tabular}
\end{solution}

\subquestion{Calculate the average marginal effect of having a parent with a 4-yr college degree. Calculate
the average marginal effect of having a parent with a 4-yr college degree for those students who
have a parent with a 4-yr college degree. Provide some intuition for why the latter marginal
effect is smaller.}

\begin{solution}
  \begin{tabular}{lcc}
    \toprule
    Variable & Average Marginal Effect \\
    \midrule
    Parent College Grad (Overall)     & 0.2449 \\
    Parent College Grad (Conditional) & 0.2257 \\
    \bottomrule
  \end{tabular}

  The average marginal effect is smaller for those who have a parent with a 4-yr degree because their baseline probability is already higher, so the effect of having another parent is smaller in terms of the change in probability.
\end{solution}

\subquestion{Use the delta method to calculate the standard error of the average marginal effect of having a parent with a 4-yr college degree. Use a non-parametric bootstrap to calculate the standard error of the average marginal effect of having a parent with a 4-yr college degree.}

\begin{solution}
  \begin{tabular}{lc}
    \toprule
    Method & Standard Error \\
    \midrule
    Delta Method       & 0.0140 \\
    Non-parametric Bootstrap & 0.0138 \\
    \bottomrule
  \end{tabular}
\end{solution}

\subquestion{Estimate a probit model of the 4-yr college choice using Xi and Z as regressors. Report the coefficient estimates and standard errors.}
\begin{solution}
  \begin{tabular}{lcc}
    \toprule
    Variable & Coefficient & Std. Error \\
    \midrule
    Constant            & -0.7069 & 0.1453 \\
    Parent College Grad & 0.7387  & 0.0448 \\
    High School GPA     & 0.3974  & 0.0513 \\
    Distance            & -0.0136 & 0.0013 \\
    \bottomrule
  \end{tabular}
\end{solution}

\subquestion{Calculate the AME of having a parent with a 4-yr college degree using the probit model. Compare with the logit model.}
\begin{solution}
  \begin{tabular}{lcc}
    \toprule
    Variable & AME (Logit) & AME (Probit) \\
    \midrule
    Parent College Grad & 0.2449 & 0.2429 \\
    \bottomrule
  \end{tabular}

  They are very similar because, in the binary choice setting, the logit and probit are often close to each other in terms of prediction. 
\end{solution}

\subquestion{Would you want to give the logit or probit coefficients or marginal effects causal interpretations? Why or why not?}
\begin{solution}
  No, we would not want to give either model a causal interpretation. In the current specification, we are not controlling for unobservables that can be correlated with the regressors. This would likely lead to a large error term, leading to biased estimates.
\end{solution}

\question{Use “dataHW2 Problem2.mat” to complete this exercise. In this problem you will work with a
multinomial choice model. Upon graduating high school, individuals have four options: work (Y=0),
attend a 2-year college (Y=1), attend a 4-year public college (Y=2), or attend a 4-year private college
(Y=3). The available student demographic characteristics (and choice specific constants) are as in
the previous example. The three columns of Xi are the constant, an indicator for whether either
parent graduated from a 4-year college, and the student’s high school GPA respectively. There are
two choice specific variables, distance (Zdist) and price (Zprice) measured in thousands of dollars.
Each column of these matrices reflects the choice specific attribute for the corresponding choice. In
other words, the first column are the attributes for the work choice, the second column for the 2-year
choice, and so on. Note that I assume that distance and price for the work option are zero.}

\subquestion{Estimate a multinomial logit model with Xi and Zdist and Zprice as regressors using the default
fminunc options. Normalize the utility of work to zero. Report parameter estimates and SEs}
\subquestion{Re-estimate the model from part (a) (again using the default fminunc options), but use different
starting values (all zeros/random draws). Report parameter estimates and SEs. Do you get
the same results? If not, what do you think is happening?}
\subquestion{Use the “optimoptions” command to set the options such that the estimation results using
either the starting values from (a) and (b) are identical.}

\begin{solution}
  So, I will begin by saying my minimize function does give me the same answers regardless of starting value but I will say why they would differ. SciPy's minimize function uses BFGS by default, a quasi-newton algorithm.
  Upon talking to classmates, this is what converged their estimates regardless of starting values. In the fminunc case, it is likely the default value gets stuck at local minimia/maxima, thus resulting in sensitivity to 
  starting value. 

  So, I kind of sidestepped this question, but I did play around with options. Any changes to tolerances just made the estimates diverge. The results can be seen in code output, and because it would be redundant to list the estimates twice, I will list the recovered estimates and std. errors.

  \begin{tabular}{lcc}
    \toprule
    Variable & Estimate & Std. Error \\
    \midrule 
    $\beta_{1,0}$ & -1.2554 & 0.1722 \\
    $\beta_{1,BA}$ & 0.1671 & 0.0553 \\
    $\beta_{1, GPA}$ & 0.1695 & 0.0646 \\
    $\beta_{2, 0}$ & -1.6212 & 0.1036 \\
    $\beta_{2, BA}$ & 1.2441 & 0.0383 \\
    $\beta_{2, GPA}$ & 0.8202 & 0.0342 \\
    $\beta_{3, 0}$ & -1.8249 & 0.2144 \\
    $\beta_{3, BA}$ & 1.1729 & 0.0568 \\
    $\beta_{3, GPA}$ & 0.5723 & 0.0647 \\
    \midrule
    $Z_{dist}$ & -0.0260 & 0.0008 \\
    $Z_{price}$ & -0.0546 & 0.0057 \\
    \bottomrule
  \end{tabular}
\end{solution}

\subquestion{Find the average marginal effect of having a parent with a 4-yr college degree on the likelihood
of attending any 4-year college}
\begin{solution}
  \begin{tabular}{lc}
    \toprule
    AME & 0.3074 \\
    \bottomrule
  \end{tabular}
\end{solution}

\subquestion{Find the average predicted probability for each choice in the sample. These numbers should
match the empirical probabilities observed in the data (check this). What parameters in the
model ensure this is the case?}

\begin{solution}
  \begin{tabular}{lc}
    \toprule
    Choice & Probability \\
    \midrule
    Work & 0.4476 \\
    2-yr & 0.1100 \\
    4-yr (public) & 0.3377 \\
    4-yr (private) & 0.1047 \\
    \bottomrule
  \end{tabular}
  
  These numbers do in fact match the data probabilities (see code output). It should be intercepts. They are the only parameter which is only variable with $j$, meaning they are the only candidate for balancing predicted and observed shares.
\end{solution}

\subquestion{Now, assume that the 4-year private option is no longer available. Find the predicted proba-
bilities of the first three choices when the 4th choice is eliminated. Note you do not have to
re-estimate the model to this. Find the percent change in the predicted choice probability
for Y={0,1,2} (relative to when 4-year private is included) and average this across the sample.
These numbers should be the same for Y={0,1,2}. What feature does this reflect? Do these
substitution patterns make sense in the context of the current problem?}

\begin{solution}
  \begin{tabular}{lcc}
    \toprule
    Choice & Probability & $\%\Delta$ \\
    \midrule
    Work & 0.4907 & 0.1260 \\
    2-yr & 0.1212 & 0.1260 \\
    4-yr (public) & 0.3880 & 0.1260 \\
    \bottomrule
  \end{tabular}

  The sameness of percent change means that people will equally spread in the probabilty distribution to other choices when an option is removed. These substitution patterns are interesting. I would have initially thought that people would fully substitute to 4 year public institutions
  upon removal of 4 year private institutions. However, with differing Z's by individual, it makes sense that the sorting to the other three would occur equally, balancing disutility of distance and price with the choices.

\end{solution}

\subquestion{Find the average change in expected consumer surplus (in thousands of dollars) when the
private 4-year option is eliminated. Find the average change in expected consumer surplus
separately for individuals whose parents do and do not have a college degree. Note that this is
a case where utility is likely not quasi-linear in income, but we’ll do the exercise for practice.}

\begin{solution}
  \begin{tabular}{lc}
    \toprule
    Case & Avg. Change in Surplus (\$) \\
    \midrule
    Remove 4-yr (all) & -2095.83 \\
    Remove 4-yr (Parent with BA) & -3872.60 \\
    Remove 4-yr (Parent without BA) & -1119.79\\
    \bottomrule
  \end{tabular}

  Most of the loss is driven by those with parents who hold a BA. This makes sense, as they are the most likely to go to college and thus would lose out more.
\end{solution}

\end{document}
